\documentclass[10pt,a4paper]{amsart}
\usepackage[margin=19mm]{geometry}
\usepackage{amsmath,amssymb,mathtools}
\usepackage[scr=rsfs,cal=euler]{mathalfa}
\usepackage{xfrac}
\usepackage[unicode,hidelinks,bookmarks=false]{hyperref}
\newcommand{\ie}{i.e.\ }
\newcommand{\cf}{cf.\ }
\newcommand{\resp}{respectively\ }
\newcommand{\Subsection}{\subsection}
\providecommand{\nobreakdash}{\nobreak\hbox{-}}
\setlength{\emergencystretch}{2em}
\multlinegap0pt
\usepackage{bm}
\usepackage{bbm}
\usepackage{dsfont}
\usepackage{xspace}
\usepackage{cancel}
\usepackage{mathtools}
\usepackage{amsthm}
\makeatletter
\@ifundefined{lemma}{\newtheorem{lemma}{Lemma}}{}
\@ifundefined{theorem}{\newtheorem{theorem}{Theorem}}{}
\makeatother
\makeatletter
\expandafter\ifx\csname keywords\endcsname\relax
\newenvironment{keywords}
  {\small\noindent\textbf{Keywords: }}
\fi
\makeatother
\title[Simple Verification and Implementation of Observer Error Dyn]{Simple Verification and Implementation of Observer Error Dynamics Linearization: A Pascal's Triangle--Hessian Matrix Criterion}
\author{Xu Haotian, Xinquan Shao, Liu Shuai}
\date{Source: arXiv:2608.18804v1 (CC BY 4.0)}
\begin{document}
\maketitle

\noindent\emph{Source and licence.} Simple Verification and Implementation of Observer Error Dynamics Linearization: A Pascal's Triangle--Hessian Matrix Criterion. Version arXiv:2608.18804v1 (CC BY 4.0), distributed under the Creative Commons Attribution 4.0 International licence, as recorded in the arXiv licence metadata for this exact version and shown on the abstract page https://arxiv.org/abs/2608.18804v1. Full rights notice: licenses/arxiv-2608.18804-CC-BY-4.0.txt.\par\smallskip

\noindent\textit{Editorial scope.} Counted unit: the Theorem whose source label is \texttt{thm:Lie\_PH\_equiv} in the article (source0.tex lines 981--983, source label \texttt{thm:Lie\_PH\_equiv}), delivered together with the article's own complete original proof of it (source0.tex lines 985--1104, one \texttt{proof} environment of 120 lines). No mathematics is written, repaired or re-derived: statement and proof are the article's own text line for line. Required context retained uncounted: sys3 (lines 158--160); observerform-trans (lines 247--249); eq:chain (lines 251--253); fms (lines 264--266); transform1 (lines 268--271); transform2 (lines 268--271); thm:PH (lines 320--327); lem-1 (lines 397--399); lem:straightening (lines 977--979). Every definition, display and label of those lines is retained unchanged. Omitted from this package and not redistributed: 1--157, 161--246, 250--250, 254--263, 267--267, 272--319, 328--396, 400--976, 980--980, 984--984, 1105--1139. Nothing inside a retained excerpt is omitted or altered: every retained body is delivered exactly as the article's lines are written, so no omission receipt of any kind accompanies this package. Citations: the retained excerpts carry no citation command at all, so no bibliography entry of the article is transcribed and no reference is redistributed. Reference adaptation: none was needed and none was made; every reference inside the delivered unit resolves inside it, so no unresolved reference or citation is produced. Printed slips kept literal: no printed word, symbol, hypothesis or number of the counted statement or of its proof was altered. Layout and class: the article's own class is \texttt{article} with packages including geometry, amsmath, amssymb, amsfonts, amsthm, bm, mathrsfs, caption, graphicx, array, multirow, booktabs, setspace, tabularx; the wrapper is the lane's pinned \texttt{amsart} editorial wrapper at 10pt on A4, which restates the theorem-like environments the retained text needs and adds the article's own macro definitions that the retained text needs (none), with extra wrapper packages (bm, bbm, dsfont, xspace, cancel, mathtools). The delivered numbering is the unit's own. Assets: no retained span contains a figure environment, a graphics-inclusion command, a PDF-page inclusion, a table or a TikZ picture; no image file is copied into this package, the package declares no asset dependency and no image file is redistributed with it. Licence: version arXiv:2608.18804v1 of this article is distributed under the Creative Commons Attribution 4.0 International licence, as recorded in the arXiv licence metadata for this exact version and shown on the abstract page https://arxiv.org/abs/2608.18804v1. A full rights notice is delivered at licenses/arxiv-2608.18804-CC-BY-4.0.txt. Characters: every retained excerpt is pure ASCII, so the delivered engine, XeLaTeX with its default Latin Modern text font, reports no missing character.
\begin{align}
	x_1^{(n)}=p(x),~y=x_1\label{sys3}
\end{align}

\begin{align}
	x_1=h(\chi),~x_2=L_fh(\chi),~\ldots,~x_n=L_f^{n-1}h(\chi)\label{observerform-trans}
\end{align} 

\begin{align}\label{eq:chain}
	\dot{x}_1 =& x_2,\quad \dot{x}_2 = x_3,\quad \ldots,\quad \dot{x}_{n-1} = x_n,\notag\\ \dot{x}_n =& p(x_1, \ldots, x_n).
\end{align}

\begin{align}
	x_1^{(n)}=\sum_{r=0}^{n-1}B_r^{(r)}(x_1),~y=x_1\label{fms}
\end{align}

\begin{align}
	&z_n=x_1,\label{transform1}\\
	&z_k=x_{n-k+1}-\sum_{r=k}^{n-1}B_r^{(r-k)}(x_1),~k=1,\ldots,n-1,\label{transform2}
\end{align}

\begin{theorem}\label{thm:PH}
	Given the high-order system~(\ref{sys3}), there exists a family of smooth univariate functions:
	\begin{align}
		B_0(x_1) &= p(x_1, 0, \ldots, 0), \label{eq:B0} \\
		B_r(x_1) &= \int \frac{\partial p}{\partial x_{r+1}}(x_1, 0, \ldots, 0) \, {\rm d}x_1,\quad r = 1, \ldots, n-1, \label{eq:Br}
	\end{align}
	such that the high-order system can possess the high-order fully measured form~(\ref{fms}) (equivalently, $p(x)=\sum_{r=0}^{n-1} D^r B_r(x_1)$) in the region $\Psi(\mathcal{U})$ (where $\Psi$ is the diffeomorphism defined in~(\ref{observerform-trans})) if and only if $p(x)$ satisfies the Pascal-Hessian condition in the region $\Psi(\mathcal{U})$. 
\end{theorem}

\begin{lemma}\label{lem-1}
	Let $F(x_1)$ be a smooth univariate function. Then for any $r=0,\ldots,n-1$, $D^r F(x_1)$ depends at most on $x_1,\ldots,x_{r+1}$.   
\end{lemma}

\begin{lemma}[Commuting Frame Rectification Theorem]\label{lem:straightening}
	Let $\tau_1,\ldots,\tau_n$ be smooth vector fields on an $n$-dimensional smooth manifold that are everywhere linearly independent and pairwise commuting (i.e., $[\tau_i,\tau_j]=0$, $\forall i,j$). Then there exists a local coordinate chart $s=(s_1,\ldots,s_n)$ such that $\tau_i=\frac{\partial}{\partial s_i}$ for each $i=1,\ldots,n$.
\end{lemma}

\begin{theorem}\label{thm:Lie_PH_equiv}
	In the observer form~(\ref{eq:chain}), $[\tau_i,\tau_j]=0$ for all $i,j=1,\ldots,n$ is equivalent to $p(x)$ satisfying the Pascal-Hessian condition.
\end{theorem}

\begin{proof}
	We prove the two directions separately.
	
	\noindent (1) Pascal-Hessian condition $\Rightarrow [\tau_i,\tau_j]=0$.

	By the sufficiency part of Theorem~\ref{thm:PH}, the Pascal-Hessian condition is equivalent to the existence of a family of smooth univariate functions $B_0(x_1),\ldots,B_{n-1}(x_1)$ such that
	\begin{equation}\label{eq:p_sum}
		p(x)=\sum_{r=0}^{n-1}D^rB_r(x_1).
	\end{equation}
	Define the coordinate transformation $x\mapsto z$ as~(\ref{transform1})--(\ref{transform2}); its Jacobian is anti-triangular with unit diagonal, hence a local diffeomorphism.

	Since $\dot{B}(x_1)=DB(x_1)$ for any $B(x_1)$, whence
	\[
	\frac{\mathrm{d}}{\mathrm{d}t}D^{r-k}B_r(x_1)=D^{r-k+1}B_r(x_1).
	\]
	Verifying the dynamics in $z$-coordinates: for $k=1$,
	\begin{align}
		\dot{z}_1=&\dot{x}_n-\sum_{r=1}^{n-1}\frac{\mathrm{d}}{\mathrm{d}t}D^{r-1}B_r(x_1)\notag\\
		=&p-\sum_{r=1}^{n-1}D^rB_r(x_1)=B_0(x_1)=B_0(z_n).
	\end{align}
	For $k\geq2$,
	\[
	\begin{aligned}
		\dot{z}_k&=x_{n-k+2}-\sum_{r=k}^{n-1}D^{r-k+1}B_r(x_1)\\
		&=z_{k-1}+B_{k-1}(z_n).
	\end{aligned}
	\]
	Therefore, the vector field $f_p$ in $z$-coordinates can be expressed as
	\begin{equation}\label{eq:fp_s}
		f_p=B_0(s_n)\frac{\partial}{\partial z_1}+\sum_{k=2}^{n}\left(z_{k-1}+B_{k-1}(z_n)\right)\frac{\partial}{\partial z_k}.
	\end{equation}

	Now compute the $\tau$ vector fields in $z$-coordinates. In $x$-coordinates, $\tau_1=e_n=\partial/\partial x_n$. From the dependence relations of the terms in~(\ref{transform1})--(\ref{transform2}) (Lemma~\ref{lem-1} guarantees that $\sum_{r=k}^{n-1}D^{r-k}B_r(x_1)$ depends only on $x_1,\ldots,x_{n-1}$), we obtain
	\[
	\tau_1s_1=1,\qquad \tau_1s_k=0\;(k\geq2).
	\]
	Hence, in $s$-coordinates, $\tau_1=\partial/\partial z_1$.

	From~(\ref{eq:fp_s}), direct computation yields
	\[
	\left[\frac{\partial}{\partial z_i},f_p\right]=\frac{\partial}{\partial z_{i+1}},\qquad i=1,\ldots,n-1.
	\]
	Furthermore, since $\tau_{i+1}=[\tau_i,f_p]$ and the Lie bracket is invariant under coordinate transformations, induction yields
	\[
	\tau_i=\frac{\partial}{\partial z_i},\qquad i=1,\ldots,n.
	\]
	Therefore, for any $i,j$,
	\[
	[\tau_i,\tau_j]=\left[\frac{\partial}{\partial z_i},\frac{\partial}{\partial z_j}\right]=0.
	\]
	The forward direction is proved.

	\medskip
	\noindent (2) $[\tau_i,\tau_j]=0\Rightarrow$ Pascal-Hessian condition.

	Assume $[\tau_i,\tau_j]=0$ for all $i,j$. By Lemma~\ref{lem:straightening}, there exist local coordinates $s=(s_1,\ldots,s_n)$ such that
	\begin{equation}\label{eq:tau_s}
		\tau_i=\frac{\partial}{\partial s_i},\qquad i=1,\ldots,n.
	\end{equation}
	From $\tau_i x_1=\partial x_1/\partial s_i$ together with $\tau_n x_1=1$ and $\tau_i x_1=0$ ($i<n$), we obtain $\partial x_1/\partial s_i=\delta_{in}$. By a suitable additive constant shift, we may take
	\begin{equation}\label{eq:sn_x1}
		s_n=x_1.
	\end{equation}

	In $s$-coordinates, express $f_p$ as
	\[
	f_p=\sum_{j=1}^n a_j(s)\frac{\partial}{\partial s_j}.
	\]
	From the recursion $\tau_{i+1}=[\tau_i,f_p]$ and~(\ref{eq:tau_s}), we obtain
	\[
	\frac{\partial}{\partial s_{i+1}}=\left[\frac{\partial}{\partial s_i},f_p\right]=\sum_{j=1}^n\frac{\partial a_j}{\partial s_i}\frac{\partial}{\partial s_j},\qquad i=1,\ldots,n-1.
	\]
	Comparing coefficients gives
	\begin{equation}\label{eq:da}
		\frac{\partial a_j}{\partial s_i}=\delta_{j,i+1},\qquad i=1,\ldots,n-1.
	\end{equation}
	Integrating~(\ref{eq:da}) successively: for $j=1$, $\partial a_1/\partial s_i=0$ ($\forall i\leq n-1$), hence $a_1=B_0(s_n)$; for $j=2$, $\partial a_2/\partial s_1=1$, $\partial a_2/\partial s_i=0$ ($i\geq2$), and integrating gives $a_2=s_1+B_1(s_n)$. Proceeding analogously, we obtain
	\begin{align}
		a_1=&B_0(s_n),\\
		a_j=&s_{j-1}+B_{j-1}(s_n),\quad j=2,\ldots,n,\label{eq:a_j}
	\end{align}
	where $B_0(s_n),\ldots,B_{n-1}(s_n)$ are smooth functions depending only on $s_n=x_1$.

	Now return from $s$-coordinates to $x$-coordinates. In the observer form~(\ref{eq:chain}), we have $x_{m+1}=L_{f_p}^m x_1$ ($m=0,\ldots,n-1$). Back-substituting step by step from~(\ref{eq:a_j}): first $x_1=s_n$. From $\dot{s}_n=s_{n-1}+B_{n-1}(x_1)=\dot{x}_1=x_2$, we obtain
	\[
	x_2=s_{n-1}+B_{n-1}(x_1).
	\]
	Continuing recursively, generally we have
	\begin{equation}\label{eq:x2s_inv}
		x_{n-k+1}=s_k+\sum_{r=k}^{n-1}L_{f_p}^{r-k}B_r(x_1),\qquad k=1,\ldots,n.
	\end{equation}
	Taking $k=1$,
	\[
	x_n=s_1+\sum_{r=1}^{n-1}L_{f_p}^{r-1}B_r(x_1).
	\]
	Applying $L_{f_p}$ to the above equality and using $\dot{s}_1=a_1=B_0(x_1)$,
	\begin{equation}\label{eq:p_Lf}
		p(x)=L_{f_p}x_n=B_0(x_1)+\sum_{r=1}^{n-1}L_{f_p}^r B_r(x_1).
	\end{equation}

	Finally, we prove that for any univariate function $B(x_1)$ and $r\leq n-1$,
	\begin{equation}\label{eq:Lf_D}
		L_{f_p}^rB(x_1)=D^rB(x_1).
	\end{equation}
	The case $r=0$ is obvious. Assume $L_{f_p}^rB(x_1)=D^rB(x_1)$ depends only on $x_1,\ldots,x_{r+1}$ (by Lemma~\ref{lem-1}). If $r\leq n-2$, then $r+1\leq n-1$, so $\partial_n\left(L_{f_p}^rB(x_1)\right)=0$. From $f_p=D+p\partial_n$, we obtain
	
	\begin{align}
		L_{f_p}^{r+1}B(x_1)&=(D+p\partial_n)\left(L_{f_p}^rB(x_1)\right)\notag\\
		&=D\left(D^rB(x_1)\right)+p\cdot0=D^{r+1}B(x_1).
	\end{align}
	Induction up to $r=n-2$ already covers all cases with $r\leq n-1$, hence~(\ref{eq:Lf_D}) holds.

	Substituting~(\ref{eq:Lf_D}) into~(\ref{eq:p_Lf}) gives
	\[
	p(x)=\sum_{r=0}^{n-1}D^rB_r(x_1).
	\]
	By Theorem~\ref{thm:PH}, the above is equivalent to the Pascal-Hessian condition. The backward direction is proved.

	In summary, Theorem~\ref{thm:Lie_PH_equiv} is established, i.e., the Pascal-Hessian condition and the Lie bracket involutivity condition $[\tau_i,\tau_j]=0$ are fully equivalent. This result demonstrates that the algebraic criterion proposed in this paper is in complete theoretical accord with the traditional differential geometric approach.
\end{proof}

\end{document}
